\section{What the two systems establish together}\label{sec:synthesis}

The two applications share no hydrology, no biology and no institution; they share a frame, a protocol and a verifier. Table~\ref{tab:resolved} lists every question the frame was asked, with the answer and the clause of Section~\ref{sec:frame} that supplies it. Three results are joint.

\begin{table}[t]
\centering
\caption{What the frozen protocols resolve, and what they do not, in the two systems.}
\label{tab:resolved}
\small
\begin{tabular}{p{0.27\textwidth}p{0.32\textwidth}p{0.32\textwidth}}
\toprule
Question & Northern cod & Edwards Aquifer \\
\midrule
Is the reference point self-viable under some constant control? & yes under q10 ($C^*=91.6$ kt); no under q05/worst ($C^*<0$) & yes at 618 ft (any cut $\ge 7.2\%$); no at 660 ft under any floor ($H^*_0=647.3$--657.9 ft) \\
Is any non-empty $T=\infty$ kernel available? & q10: for $C\le C_{\mathrm{vac}}=215.2$ kt; q05: $C\le 8.7$ kt; worst: none & 618 ft: yes for every cut $\ge 7.2\%$; 660 ft: none \\
Does a staged rule beat the flat cap it coincides with at the threshold? & no (S1, cascade $\equiv$ flat 60 kt on the safe set) & no in protection; yes in supply (S1 $+3.3\%$ over flat 90\%, CPM $+50.6\%$ over flat 60\%) \\
Horizon certified under the declared defect (exact kernel) & 8--13 yr (worst/q05), 12--28 yr (q10); Lipschitz: 5--7 yr & 3 yr positive-pumping (4 zero); 2 yr OOS positive-pumping (3 zero) \\
Is the horizon policy-sensitive? & exact: yes (zero catch 12 vs 240 kt 8 yr, worst floor); Lipschitz: no & weakly (3 vs 4 yr) \\
What sets the horizon? & local expansion at the LRP, $F'(\theta)=1.153$ & defect-to-contraction ratio $\varepsilon/(1-a)=60.7$ ft against a 92-ft band to the ceiling \\
Is the key constant identified? & $g(\theta)$ and $C^*$ yes (profile), $K$ no & $a$, $H^*$ yes (bootstrap SD 10--13 ft on $H^*$) \\
Stochastic (non-worst-case) reading & 20-yr survival from the LRP 0.90 (BAU) to 0.65 (120 kt) & $P(H^*\ge618)$: 0.41 (BAU), 0.53 (flat 90\%), 0.65 (S1), 0.99 (zero) \\
\bottomrule
\end{tabular}
\end{table}


\paragraph{Protection and supply are one budget at the threshold.}
At a threshold $\theta$ the kernel of a constant-control rule is decided by the sign of $F(\theta,u)-\theta$ under the worst floor; by Proposition~\ref{prop:stage} a staged rule is decided by its stage at $\theta$. Every rule therefore has one number at $\theta$ and a one-dimensional margin, $\Cstar-C(\theta)$ for the stock and $\Hstar(P(\theta))-\theta$ for the aquifer. A reactive rule cannot out-supply, at the threshold, the flat rule it matches there; where it can out-supply is away from the threshold, and that is where the retention rule looks. For cod the Stage-I rule supplies $10.0\kt$ on the observed states against the flat 60-kt cap's 60, and the two are equally protective: it fails clause (b) and is not retained. For the aquifer the Stage-I rule supplies $3.3\%$ more than flat 90\% on the observed heads at equal nominal protection and $16.2\%$ more than flat 80\% at equal certified protection, and is retained. The rule is the same; the verdicts differ because the observed states of the two systems sit on different sides of the stage boundaries.

\paragraph{Certification has a horizon, and the map sets it.}
The certified horizon of the two systems is set by different features of the fitted map and the declared defect. The cod map is expansive at the LRP ($F'(\theta)=1.153$ under the fitted $K$), so the Lipschitz error grows geometrically and its horizon is six years under the worst/q05 floors and seven under q10 for catches up to 150 kt; the exact bounded-defect kernel extends this to 8--13 years under the two harsher floors and 12--28 under q10, with catch-sensitive horizons (Figure~\ref{fig:horizons}). The aquifer map is contractive ($a=0.746$) and the error saturates at $r_\infty=60.7\ft$. At 618 ft the declared-defect horizon is three years for positive-pumping rules and four at zero pumping; with the out-of-sample defect the horizons are two and three. The exact kernel matches the Lipschitz boundary for positive constant-pumping rules wherever the backward recursion binds, while zero pumping at $T=1,2$ and the staged rules depart locally without changing a horizon (Proposition~\ref{prop:lip}(ii)). Thus the cod Lipschitz result is conservative relative to the exact recursion, whereas the aquifer's positive-pumping horizon is class-limited under the frozen protocol. The mechanisms of boundary departure also differ: curvature lowers the cod boundary away from the threshold; in the aquifer, the inflated initial-state constraint binds at zero pumping and stage crossings make the staged maps piecewise affine.

\begin{figure}[t]\centering
\includegraphics[width=\textwidth]{fig_horizons.pdf}
\caption{Certified horizons. (a)~Cod LRP: last non-empty horizon of the exact bounded-defect kernel under the three floors (bars) and of the Lipschitz certificate (black ticks; for the staged rules the certificate of their 60-kt stage at the LRP). (b)~Edwards 618-ft trigger under the 1956 floor: last non-empty horizon of the exact bounded-defect kernel with the declared defect ($15.4\ft$) and with the out-of-sample defect ($21.8\ft$); the Lipschitz certificate coincides with the exact kernel wherever the backward recursion binds (Proposition~\ref{prop:lip}(ii)); the zero-pumping and staged-rule boundary departures are described in Section~\ref{sec:edw-horizon} and do not change the horizons.}
\label{fig:horizons}
\end{figure}


\paragraph{Thresholds can be unsecurable, and the frame says which.}
Theorem~\ref{thm:margin} reduces securability to one comparison per floor. The cod LRP is secured by any constant catch below $\Cstar$ under every floor; the Edwards 618-ft trigger is secured by any cut of at least $7.2\%$ under the worst floor and by BAU under the two quantile floors; the Edwards 660-ft trigger is secured by nothing, under any declared floor, because the zero-pumping attractor falls short of it by $2.1$--$12.7\ft$. A test that asks only whether a threshold is breached in simulation reports the 660-ft line as frequently crossed and recovered; the frame reports that no admissible rule can be credited with its recovery, which is the operational distinction between a line that is managed \emph{to} and a line that is managed \emph{around}.

\paragraph{Cadence.}
The certified horizons bound the re-certification interval. With a three-year horizon on the declared defect and two on the audited one, the aquifer's registered map should be refitted and the kernels recomputed at least every two years; the cod exact horizon of eight to thirteen years under the harsher floors, read against the six-year review cycle of the International Whaling Commission's management procedures and the one-to-three-year stock-assessment cadence of most North American groundfish, says that the present cadence is inside the certificate for cod and outside it for the aquifer under BAU. Both protocols were frozen on 2026-08-26, before any kernel, boundary, replay or retention score was computed, and both report negative results; a frame that returned only positive verdicts would be advocacy.
